\section{Generality and limits of the test}
\label{sec:generality}

The preceding sections rest on one dataset, resolution and task.
Section~\ref{sec:dreambooth} changes the model, resolution and task together,
in a setting where the metric precondition P4 fails;
Section~\ref{sec:celeba} moves to another dataset and resolution, where the
price of the same spectral correction changes sign and depends on the
operator. Gaussian denoising, which was not priced and shows no difference
from MSE after correction, is reported in Appendix~\ref{sec:restoration}.

\subsection{Subject-driven fine-tuning}
\label{sec:dreambooth}

We fine-tune the U-Net of Stable Diffusion~1.5 \citep{rombach2022high} with
DreamBooth \citep{ruiz2023dreambooth} on one subject (five instance
photographs) for 400 steps at $512\times512$, with \dbNobjectives{}
objectives and five seeds each; AWWL uses the conference paper's setting for
this task, $\alpha = 0.8$, $p = 2$ (Section~\ref{sec:family-awwl}). Each run
generates \dbNimages{} images from \dbNprompts{} prompts placing the subject
in new scenes or styles, scored with CLIP ViT-B/32
\citep{radford2021learning} on text--image agreement and on similarity to the
instance photographs (Appendix~\ref{app:setup-dreambooth}); DINO similarity
\citep{caron2021emerging} was not computed (Section~\ref{sec:limitations}).

\guardedfigure{0.64}{dreambooth_tradeoff.png}{%
  DreamBooth on Stable Diffusion~1.5 (one subject, five seeds): paired
  difference from MSE (the dot at the origin) in CLIP text--image score and
  subject similarity, with paired 95\% $t$-intervals. Filled: differs from MSE
  after Holm correction on at least one metric (only raw SNR weighting, on
  CLIP score). None lies in the upper-right quadrant (better than MSE on both
  metrics). Means and ranks: Table~\ref{tab:dreambooth}.}{fig:db-tradeoff}

The conference paper \citep{thao2026awwl} reports AWWL as the only method in
the top two on both metrics. Over five seeds it is ranked \dbRankSimAWWL{} of
\dbNobjectives{} on similarity but \dbRankClipAWWL{} on CLIP score, below MSE
(Table~\ref{tab:dreambooth}), and neither difference from MSE survives Holm
correction (Fig.~\ref{fig:db-tradeoff}; $p_{\mathrm{Holm}} = \dbDSimAWWLHolm$
and \dbDClipAWWLHolm). The only surviving difference is the gain of raw SNR
weighting in CLIP score ($+\dbDClipRawSNR$,
$p_{\mathrm{Holm}} = \dbDClipRawSNRHolm$), an arm that is last on similarity
and fails to train from scratch on CIFAR-10 (FID
\num[round-mode=places,round-precision=0]{\fidRawSNRmean};
Section~\ref{sec:compare}). Neither Min-SNR-$\gamma$ nor any
frequency-aware objective differs significantly from MSE on this task.

\paragraph{Pricing.} AWWL changes the high band here, but the metrics cannot
price the change. Against the instance photographs, MSE's samples lie
\dbDefMSE~dB below the reference in the high band and AWWL's \dbDefAWWL~dB, a
paired correction of \dbCorr~dB $[\dbCorrLo, \dbCorrHi]$ (raw
$p = \dbCorrP$); as the reference lacks the prompted scenes and styles, this
deficit mixes content mismatch with spectral bias. Boosting each run's high
band by its full deficit (leaving MSE's samples \dbResid~dB below the
reference) changes CLIP score and similarity by means below $10^{-4}$ in
magnitude, with standard deviations of
\num[exponent-mode=scientific]{\dbPriceClipSD} and
\num[exponent-mode=scientific]{\dbPriceSimSD} over the \dbPriceClipN{} MSE
and AWWL runs, against seed standard deviations of \dbSeedSDClip{} and
\dbSeedSDSim{}. This price is indistinguishable from zero, although the
boosted images were stored as dithered 8-bit files (contrary to P0), and
follows from the metric's design: the boost acts only above two thirds of the
Nyquist frequency of the 512-pixel samples, whereas the CLIP encoder receives
an antialiased 224-pixel resize, which attenuates everything above $224/512$
of that frequency.
Precondition P4 (Section~\ref{sec:pricing-pre}) therefore fails, which
limits the test rather than the objective: no CLIP difference on this task,
AWWL's lead in similarity included, can be credited to the high-band
correction that AWWL names as its mechanism.

\subsection{Another dataset and resolution: CelebA-64}
\label{sec:celeba}

CelebA \citep{liu2015deep} at $64\times64$ keeps the architecture family,
recipe (200 epochs), sampler and FID protocol of CIFAR-10
(Appendix~\ref{app:setup-celeba}), with FID computed against the full training
split. Five MSE seeds reach FID $\celebaFIDMSEmean \pm \celebaFIDMSEsd$. AWWL
was not trained here (Section~\ref{sec:limitations}), so only the value side
of the test is computed.

\paragraph{Coherent corrections raise FID\@.} The baseline's high-band deficit
matches CIFAR-10's (\celebaDeficitHigh~dB against \deficitHighMSE~dB), yet
correcting it makes FID worse (Table~\ref{tab:price-curve}; a negative value
means a higher FID). The $+0.41$~dB boost, which reduces the deficit by the
same amount on both datasets, is worth $+\priceBoostValue$~FID
$[\priceBoostValueLo, \priceBoostValueHi]$ on CIFAR-10 and
$\celebaPriceBoost$ $[\celebaPriceBoostLo, \celebaPriceBoostHi]$ on
CelebA-64; a $+1.25$~dB boost is worth $\celebaPriceBoostHigh$
$[\celebaPriceBoostHighLo, \celebaPriceBoostHighHi]$; and the full
radial-spectrum match, worth $+\priceMatchValue$ on CIFAR-10, is worth
$\celebaPriceMatch$ $[\celebaPriceMatchLo, \celebaPriceMatchHi]$ (raw
$p \celebaPriceMatchP$ for each). P4 holds, since FID responds more strongly
to a larger correction; what differs is the sign.

\paragraph{Incoherent noise lowers it.} Uniform white noise of $\pm0.5$~LSB
per channel, added to the same samples without rounding, is worth
$+\celebaNoiseHalfValue$ $[\celebaNoiseHalfValueLo, \celebaNoiseHalfValueHi]$,
and noise of $\pm1$~LSB $+\celebaNoiseOneValue$
$[\celebaNoiseOneValueLo, \celebaNoiseOneValueHi]$ (raw
$p \celebaNoiseOneValueP$ for each), as on CIFAR-10
(Section~\ref{sec:worth-curve}). Relative to the image spectrum, the noise,
like the boost, adds energy mostly in the high band, yet the two move FID in
opposite directions: direct evidence that $V$ depends on the operator, not
only on the statistic it changes (P3). Hazard H1 thus also arises at a second
resolution and operating point.

\paragraph{Attribution.} The two settings differ in the data; in the
resolution, and with it the factor by which FID's Inception network upsamples
the samples to 299 pixels (4.7 against 9.3); in the operating point (FID
\celebaFIDMSEmean{} against
\num[round-mode=places,round-precision=2]{\fidMSEmean}); and, since 200
epochs cover a larger training set, in the number of optimiser steps (3.25
times as many). The experiment does not separate these factors, and the
band attribution of Section~\ref{sec:worth} was not repeated on CelebA-64.
The price of the same spectral correction thus differs in sign between two
settings, and a mechanism that pays in one cannot be assumed to pay in
another. Because the value side needs only scoring passes, it can be computed
before any training; on CelebA-64 it predicts that an objective whose only
effect was this coherent correction would raise FID.
